\section{Архитектура веб-приложения и стек вызовов}
\label{sec:architecture}

Веб-приложение построено по многоуровневой схеме. Прохождение любого пользовательского действия организовано в виде строгого сквозного стека вызовов:
\begin{enumerate}
    \item клиентский уровень: веб-страница браузера (HTML-формы с параметрами операций DDL, DML или DQL);
    \item уровень веб-сервера и маршрутизации: маршрутизатор стандартной библиотеки Golang (версия 1.27), диспетчеризация по HTTP-методу и пути, распаковка параметров формы;
    \item уровень сборщика запросов (Query Builder): экранирование идентификаторов PostgreSQL двойными кавычками (\texttt{quoteIdent}) и строковых литералов одинарными кавычками (\texttt{quoteLiteral});
    \item уровень СУБД PostgreSQL: выполнение запроса через пул соединений \texttt{pgx}, опрос системного каталога \texttt{information\_schema};
    \item уровень представления (View Renderer): динамическое сканирование кортежей и отрисовка HTML-шаблона с отображением текста выполненного SQL-запроса и полученных данных.
\end{enumerate}

\subsection{Веб-сервер и маршрутизация запросов}

Инициализация сетевого мультиплексора и регистрация маршрутов выполняются в функции \texttt{main} (листинг~\ref{lst:main-router}).

\begin{listing}[H]
    \gofile[firstline=70, lastline=105, firstnumber=70]{../main.go}
    \caption{Маршрутизация HTTP-запросов DDL, DML и DQL в \texttt{main.go}}
    \label{lst:main-router}
\end{listing}

В строке 73 шаблон \texttt{"GET /\{\$\}"} задает строгое соответствие корневому пути и выполняет перенаправление на страницу таблиц \texttt{/tables}. Для остальных операций предусмотрено четкое разделение:
\begin{itemize}
    \item \texttt{GET}-запросы отображают страницы и формы (например, \texttt{/ddl/\{op\}}, \texttt{/dml/\{op\}}, \texttt{/select});
    \item \texttt{POST}-запросы принимают данные форм и запускают сборку и выполнение соответствующих SQL-команд (\texttt{/ddl/create}, \texttt{/dml/insert}, \texttt{/query/select} и др.).
\end{itemize}

\subsection{Архитектура обработчиков и обработка ошибок}

Для унификации обработки ошибок используется тип-адаптер \texttt{AppHandler} (листинг~\ref{lst:app-handler}), представляющий собой функцию, возвращающую ошибку.

\begin{listing}[H]
    \gofile[firstline=131, lastline=173, firstnumber=131]{../views.go}
    \caption{Структура \texttt{AppError} и адаптер обработчиков \texttt{AppHandler}}
    \label{lst:app-handler}
\end{listing}

Метод \texttt{ServeHTTP} перехватывает ошибки выполнения. Если возвращена ошибка \texttt{AppError}, фиксируется запись в логе и клиенту отправляется соответствующий HTTP-статус. При возникновении ошибки на уровне СУБД форма рендерится повторно со статусом \texttt{400 Bad Request}, сохраняя введенные пользователем значения и подсвечивая сообщение сервера PostgreSQL и проблемный SQL-запрос.

\subsection{Движок представления и шаблонизация}

Генерация веб-интерфейса реализована на стороне сервера с помощью встроенного пакета \texttt{html/template} и обертки \texttt{ViewRenderer} (листинг~\ref{lst:view-renderer}).

\begin{listing}[H]
    \gofile[firstline=14, lastline=39, firstnumber=14]{../views.go}
    \caption{Инициализация и рендеринг шаблонов в \texttt{views.go}}
    \label{lst:view-renderer}
\end{listing}

При запуске сервера компилируется базовый шаблон \texttt{layout.html}, после чего методом \texttt{Clone()} создаются изолированные копии для каждой рабочей страницы (\texttt{tables}, \texttt{ddl}, \texttt{dml}, \texttt{select}, \texttt{join}). Для динамических вычислений в разметке зарегистрированы специальные функции (листинг~\ref{lst:func-map}).

\begin{listing}[H]
    \gofile[firstline=44, lastline=60, firstnumber=44]{../main.go}
    \caption{Пользовательские функции шаблонизатора \texttt{funcMap}}
    \label{lst:func-map}
\end{listing}

Функции \texttt{seq} и \texttt{add} используются для циклической генерации динамических строк формы (например, добавление заданного количества колонок в DDL или кортежей в DML). Функция \texttt{getRowVal} обеспечивает безопасное извлечение сохраненных значений полей при повторной отрисовке после ошибок.

Базовый шаблон \texttt{layout.html} (листинг~\ref{lst:layout-html}) содержит навигацию, кнопку быстрого сброса БД и блок отображения сгенерированного SQL-запроса. Для предотвращения выхода широких таблиц за границы страницы контейнерам \texttt{figure} назначен стиль горизонтальной прокрутки \texttt{overflow-x: auto}.

\begin{listing}[H]
    \htmlfile[firstline=16, lastline=46, firstnumber=16]{../templates/layout.html}
    \caption{Разметка контейнера и блок аудита SQL в \texttt{layout.html}}
    \label{lst:layout-html}
\end{listing}

В строках 34--41 размещен блок уведомления об ошибках: если сервер СУБД возвращает ошибку, точный текст сформированного запроса выводится в теге \texttt{<pre><code>}.

\subsection{Пул соединений и взаимодействие с PostgreSQL}

Подключение к СУБД осуществляется в функции \texttt{main} через пул соединений библиотеки \texttt{pgx} (листинг~\ref{lst:db-pool}).

\begin{listing}[H]
    \gofile[firstline=25, lastline=37, firstnumber=25]{../main.go}
    \caption{Инициализация пула соединений \texttt{pgx}}
    \label{lst:db-pool}
\end{listing}

Пул соединений повторно использует открытые TCP-сессии к PostgreSQL, снижая накладные расходы на постоянную аутентификацию и установление соединения при каждом HTTP-запросе.

\subsection{Механизм безопасной сборки SQL-запросов}

Поскольку имена таблиц и столбцов передаются из веб-интерфейса динамически, прямое использование позиционных плейсхолдеров PostgreSQL (\texttt{\$1, \$2}) для них неприменимо. Для защиты от внедрения вредоносного SQL-кода (SQL Injection) в файле \texttt{db.go} реализованы специализированные функции экранирования (листинг~\ref{lst:quote-helpers}).

\begin{listing}[H]
    \gofile[firstline=39, lastline=45, firstnumber=39]{../db.go}
    \caption{Функции экранирования идентификаторов и литералов}
    \label{lst:quote-helpers}
\end{listing}

Функция \texttt{quoteIdent} экранирует идентификаторы PostgreSQL (имена таблиц, колонок) согласно стандарту: имя оборачивается в двойные кавычки, а внутренние кавычки удваиваются (\texttt{""}). Функция \texttt{quoteLiteral} обрамляет пользовательские строковые константы в одинарные кавычки с удвоением внутренних одинарных кавычек.
